%% 19장 가우스 소거와 LU
\chapter{Gaussian Elimination and LU Factorization}
\label{ch:lu}

\abstract{연립일차방정식 $Ax=b$를 푸는 가우스 소거는 수치선형대수의 출발점이다. 이 장에서는 가우스 소거를 행렬의 분해 $A=LU$로 다시 읽고, 그 분해가 존재하는 조건과 행을 바꾸는 피벗이 왜 필요한지를 본다. 이어서 9장의 반올림 오차 해석을 가우스 소거에 적용해, 계산한 인수의 곱이 원래 행렬에 아주 가까운 행렬과 정확히 같다는 후진 오차 정리를 증명한다. 그 오차의 크기를 정하는 것은 소거하는 동안 원소가 얼마나 커지는가, 곧 성장 인수이다. 부분 피벗의 성장 인수가 최악의 경우 $2^{n-1}$까지 커질 수 있다는 것과, 그런데도 실제로는 거의 언제나 작다는 역설을 직접 계산으로 확인한다. 마지막으로 같은 알고리즘을 행렬 곱셈 위주로 다시 짜서 13장의 캐시를 잘 쓰게 하는 블록 LU와, 그 바탕인 BLAS의 층을 다룬다.}

\section{LU Factorization}
\label{sec:lu-factorization}

\subsection{Elimination as Factorization}

$n\times n$ 행렬 $A$로 $Ax=b$를 푸는 가우스 소거는 첫째 행의 배수를 아래 행들에서 빼서 첫째 열의 대각 아래를 $0$으로 만들고, 같은 일을 둘째 열, 셋째 열에 차례로 되풀이해 위삼각 행렬을 얻은 뒤, 아래에서부터 거꾸로 대입한다. $k$번째 단계에서 $i$번째 행($i>k$)에 빼는 $k$번째 행의 배수를 $l_{ik}=a_{ik}^{(k)}/a_{kk}^{(k)}$라 하자. 여기서 $a^{(k)}$는 $k$번째 단계 직전의 행렬이다. 이 단계는 단위 하삼각 행렬
\[
M_k=I-m_ke_k^{\mathsf T},\qquad m_k=(0,\dots,0,l_{k+1,k},\dots,l_{nk})^{\mathsf T}
\]
을 왼쪽에 곱하는 것과 같고, $M_k^{-1}=I+m_ke_k^{\mathsf T}$이다. 소거가 끝나면 $M_{n-1}\cdots M_1A=U$가 위삼각이므로
\[
A=M_1^{-1}\cdots M_{n-1}^{-1}U=LU,\qquad L=I+\sum_{k=1}^{n-1}m_ke_k^{\mathsf T}
\]
이다. 곱 $M_1^{-1}\cdots M_{n-1}^{-1}$이 각 $m_k$를 제자리에 그대로 모은 단위 하삼각 행렬이 되는 것은 $e_k^{\mathsf T}m_j=0$($j\ge k$)이어서 교차항이 모두 사라지기 때문이다. 곧 가우스 소거는 $A$를 단위 하삼각 행렬 $L$과 위삼각 행렬 $U$의 곱으로 분해하는 것이고, 소거의 배수들이 $L$의 대각 아래에 그대로 놓인다. 분해를 한 번 해 두면 우변이 여러 개여도 $Ly=b$와 $Ux=y$의 두 삼각 풀기로 각각 $O(n^2)$에 풀 수 있다.

연산의 수를 세어 보자. $k$번째 단계는 $(n-k)\times(n-k)$ 블록의 각 원소에서 곱 하나를 빼므로 곱셈과 뺄셈이 $(n-k)^2$번씩이다. 따라서 분해 전체의 부동소수점 연산은 $\sum_{k=1}^{n-1}2(n-k)^2\approx2n^3/3$번이다. 삼각 풀기 둘은 $2n^2$번이므로, $n$이 크면 비용의 거의 전부가 분해에 있다.

\subsection{Existence and Uniqueness}

소거는 피벗 $a_{kk}^{(k)}$가 $0$이 아닐 때만 진행된다. 이 조건을 $A$로 표현하면 다음과 같다.

\begin{theorem}
\label{thm:lu-exists}
$A$의 앞쪽 주소행렬 $A_k=A(1{:}k,1{:}k)$가 $k=1,\dots,n-1$에서 모두 정칙이면, 단위 하삼각 $L$과 위삼각 $U$로 $A=LU$인 분해가 하나뿐이다.
\end{theorem}

\begin{proof}
$n$에 대한 귀납법이다. $n=1$이면 $L=1$, $U=A$이다. $A$를
\[
A=\begin{pmatrix}A_{n-1}&c\\r^{\mathsf T}&\alpha\end{pmatrix},\qquad
L=\begin{pmatrix}L'&0\\l^{\mathsf T}&1\end{pmatrix},\qquad
U=\begin{pmatrix}U'&u\\0&\upsilon\end{pmatrix}
\]
로 나누면 $A=LU$는 $A_{n-1}=L'U'$, $c=L'u$, $r^{\mathsf T}=l^{\mathsf T}U'$, $\alpha=l^{\mathsf T}u+\upsilon$과 동치이다. $A_{n-1}$의 앞쪽 주소행렬은 $A$의 것과 같으므로 귀납 가정에서 $A_{n-1}=L'U'$가 하나뿐이고, $A_{n-1}$이 정칙이므로 $U'$도 정칙이다. 그러면 $u=L'^{-1}c$, $l^{\mathsf T}=r^{\mathsf T}U'^{-1}$, $\upsilon=\alpha-l^{\mathsf T}u$가 하나씩으로 정해진다.
\end{proof}

이 조건은 정칙인 행렬에서도 쉽게 깨진다. $\begin{psmallmatrix}0&1\\1&0\end{psmallmatrix}$은 정칙이지만 $a_{11}=0$이어서 LU 분해가 없다. 행의 순서를 바꾸면 이 문제가 사라진다.

\section{Pivoting}
\label{sec:pivoting}

\subsection{The Danger of Small Pivots}

피벗이 정확히 $0$이 아니어도 아주 작으면 문제가 된다.
\[
A=\begin{pmatrix}10^{-20}&1\\1&1\end{pmatrix},\qquad b=\begin{pmatrix}1\\2\end{pmatrix}
\]
의 해는 거의 $(1,1)$이고 $A$의 조건수는 약 $2.6$으로 아주 작다. 그런데 피벗 없이 소거하면 배수가 $l_{21}=10^{20}$이고 $u_{22}=1-10^{20}$은 binary64로 $-10^{20}$으로 반올림되어 $1$이 사라진다. 거꾸로 대입하면 $x_2=1$은 맞지만 $x_1=(1-x_2)/10^{-20}=0$이 나온다. 조건이 좋은 문제를 불안정한 알고리즘으로 풀어 답을 완전히 잃은 것이다. 원인은 작은 피벗이 만든 큰 배수가 $u_{22}$의 계산에서 원래의 정보 $a_{22}=1$을 반올림으로 지워 버렸다는 데 있다.

\subsection{Partial Pivoting}

해법은 $k$번째 단계에서 $k$번째 열의 대각 아래 원소 가운데 절댓값이 가장 큰 것을 찾아 그 행을 $k$번째 행과 바꾼 뒤 소거하는 것이다. 이것을 \term{부분 피벗}{partial pivoting}이라 한다. 그러면 모든 배수가 $|l_{ik}|\le1$을 만족한다. 행을 바꾸는 것은 왼쪽에 치환 행렬을 곱하는 것이므로, 부분 피벗 가우스 소거는 치환 행렬 $P$에 대해 $PA=LU$를 구하는 것이다. $A$가 정칙이면 각 단계의 열에 $0$이 아닌 원소가 반드시 있으므로 이 분해는 언제나 존재한다. 행과 열을 모두 바꾸어 남은 부분 행렬 전체에서 가장 큰 원소를 피벗으로 삼는 \term{완전 피벗}{complete pivoting}도 있지만, 찾는 비용이 $O(n^3)$이 되어 거의 쓰이지 않는다.

\begin{algorithm}[t]
\begin{algorithmic}[1]
\Require $n\times n$ 행렬 $A$
\Ensure $A$의 자리에 $L$(대각 아래)과 $U$(대각과 그 위), 치환 $p$, $PA=LU$
\State $p\gets(1,2,\dots,n)$
\For{$k\gets1$ \textbf{to} $n-1$}
  \State $r\gets\arg\max_{i\ge k}|a_{ik}|$ \Comment{부분 피벗}
  \State 행 $k$와 행 $r$을 바꾸고 $p_k\leftrightarrow p_r$
  \For{$i\gets k+1$ \textbf{to} $n$}
    \State $a_{ik}\gets a_{ik}/a_{kk}$ \Comment{배수 $l_{ik}$, $|l_{ik}|\le1$}
  \EndFor
  \For{$i\gets k+1$ \textbf{to} $n$, $j\gets k+1$ \textbf{to} $n$}
    \State $a_{ij}\gets a_{ij}-a_{ik}a_{kj}$ \Comment{계수 하나 줄인 갱신}
  \EndFor
\EndFor
\end{algorithmic}
\caption{부분 피벗 LU 분해}
\label{alg:lu-pp}
\end{algorithm}

\section{Rounding Error Analysis}
\label{sec:lu-error}

\subsection{Backward Error}

9장의 표준 모형으로 \cref{alg:lu-pp}을 분석하자. 피벗으로 행을 바꾸는 일에서는 오차가 생기지 않으므로 $P=I$로 두어도 된다. 소거가 끝났을 때 $u_{kj}$($k\le j$)와 $l_{ik}$($i>k$)는 각각
\[
u_{kj}=a_{kj}-\sum_{i<k}l_{ki}u_{ij},\qquad l_{ik}=\Bigl(a_{ik}-\sum_{j<k}l_{ij}u_{jk}\Bigr)\Bigm/u_{kk}
\]
를 계산한 것과 같다. 갱신을 되풀이하는 순서가 이 합을 차례로 빼는 순서와 같기 때문이다. 이 식들은 모두 9장에서 본 내적에서 한 항을 구하는 꼴이다.

\begin{theorem}
\label{thm:lu-backward}
부분 피벗 LU 분해가 계산한 인수 $\hat L$, $\hat U$는
\[
\hat L\hat U=PA+\Delta A,\qquad|\Delta A|\le\gamma_n|\hat L||\hat U|
\]
를 만족한다. 부등식과 절댓값은 원소별이다.
\end{theorem}

\begin{proof}
$P=I$로 두자. $k\le j$일 때 $\hat u_{kj}=\fl\bigl(a_{kj}-\sum_{i<k}\hat l_{ki}\hat u_{ij}\bigr)$이다. 이 값을 $k-1$번의 곱셈과 뺄셈으로 계산하면 9장의 내적 분석과 같은 방법으로
\[
\Bigl|a_{kj}-\sum_{i<k}\hat l_{ki}\hat u_{ij}-\hat u_{kj}\Bigr|\le\gamma_k\sum_{i\le k}|\hat l_{ki}||\hat u_{ij}|
\]
를 얻는다. 여기서 $\hat l_{kk}=1$로 두었다. 곧 $(\hat L\hat U)_{kj}$와 $a_{kj}$의 차가 $\gamma_k(|\hat L||\hat U|)_{kj}$ 이하이다. $i>k$인 $\hat l_{ik}$는 같은 합을 계산한 뒤 $\hat u_{kk}$로 한 번 더 나눈 것이므로, 나눗셈의 반올림 인수를 양변에 곱해 옮기면 같은 꼴의 부등식이 $\gamma_k$로 성립한다. 모든 원소에서 $\gamma_k\le\gamma_n$이므로 정리가 성립한다.
\end{proof}

삼각 풀기도 같은 방법으로 후진 안정임을 보일 수 있다. $(\hat L+\Delta L)\hat y=b$, $|\Delta L|\le\gamma_n|\hat L|$이고 $U$에 대해서도 같다\cite{higham2002accuracy}. 세 결과를 합치면 계산한 해 $\hat x$가
\[
(A+\Delta A)\hat x=b,\qquad|\Delta A|\le\gamma_{3n}|\hat L||\hat U|+O(u^2)
\]
를 만족한다. 곧 가우스 소거는 $|\hat L||\hat U|$가 $|A|$에 비해 크지 않는 한 후진 안정이다\cite{wilkinson1961error}.

\subsection{The Growth Factor}

부분 피벗에서는 $|l_{ij}|\le1$이므로 $|\hat L||\hat U|$의 크기는 $\hat U$의 원소가 얼마나 큰지에 달려 있다. 소거하는 동안 원소가 얼마나 자라는지를
\[
\rho_n=\frac{\max_{i,j,k}|a_{ij}^{(k)}|}{\max_{i,j}|a_{ij}|}
\]
로 재고 \term{성장 인수}{growth factor}라 한다. $\|\cdot\|_\infty$로 정리하면 $\||\hat L||\hat U|\|_\infty\le n^2\rho_n\|A\|_\infty$ 정도이므로, 후진 오차는 대략 $n^3\rho_nu$로 묶인다.

\begin{proposition}
\label{prop:growth-bound}
부분 피벗에서 $\rho_n\le2^{n-1}$이다.
\end{proposition}

\begin{proof}
한 단계의 갱신은 $a_{ij}^{(k+1)}=a_{ij}^{(k)}-l_{ik}a_{kj}^{(k)}$이고 $|l_{ik}|\le1$이므로 $\max|a^{(k+1)}|\le2\max|a^{(k)}|$이다. 단계는 $n-1$번이므로 $\rho_n\le2^{n-1}$이다.
\end{proof}

이 상한은 실제로 도달한다. 대각이 $1$, 대각 아래가 모두 $-1$, 마지막 열이 모두 $1$이고 나머지가 $0$인 행렬은 부분 피벗이 행을 전혀 바꾸지 않고, 소거할 때마다 마지막 열의 원소가 두 배가 되어 $\rho_n=2^{n-1}$이다. 직접 계산하면 $n=10$에서 $512$, $n=60$에서 $2^{59}\approx5.8\times10^{17}$이다. 이 행렬은 $n=60$에서 조건수가 약 $27$에 불과한데도, 무작위 해 $x$로 만든 우변을 부분 피벗 LU로 풀면 상대오차가 약 $0.59$로 답을 거의 다 잃는다. 후진 안정성이 무너졌다.

그런데 실제 응용에서 이런 일은 놀랄 만큼 드물다. 원소가 표준 정규분포인 무작위 행렬에서 부분 피벗의 성장 인수를 재어 보면 $n=10$에서 평균 $1.7$, $n=100$에서 $5.4$, $n=300$에서 $10.1$이었다. 상한 $2^{n-1}$과는 비교할 수 없이 작다. 왜 그런지에 대한 완전한 설명은 아직 없으며, Trefethen과 Schreiber는 큰 성장을 낳는 행렬들이 행렬 공간에서 지수적으로 드물다는 것을 실험과 확률적 논증으로 보였다\cite{trefethen1990average}. 미분방정식의 경계값 문제처럼 실제 응용에서 큰 성장이 생기는 예도 보고되어 있지만\cite{wright1993}, 부분 피벗 가우스 소거는 그런 예외를 감수하고도 표준 알고리즘으로 남아 있다.

\begin{aside}[가우스 소거의 계보]
방정식을 하나씩 소거해 푸는 방법은 기원전 수 세기의 중국 수학서 『구장산술』의 ``방정'' 장에 이미 셈대로 계산하는 절차로 나온다. 유럽에서는 Newton의 강의 노트가 이 방법을 대수 교과서에 퍼뜨렸고, Gauss는 소행성 궤도의 최소제곱 계산에서 정규방정식을 이 방법으로 푸는 체계적인 표기를 만들었다. 이 방법에 Gauss의 이름이 붙은 것은 20세기에 들어서이다. 행렬의 분해 $A=LU$라는 관점은 1940년대에 컴퓨터로 큰 연립방정식을 풀게 되면서 von Neumann, Turing, Wilkinson의 오차 해석과 함께 자리 잡았다\cite{grcar2011}.
\end{aside}

\section{Block Algorithms}
\label{sec:blocked-lu}

\subsection{The Levels of BLAS}

\cref{alg:lu-pp}은 갱신 단계에서 $(n-k)^2$개의 원소를 하나씩 읽고 써서, 원소 하나를 메모리에서 가져올 때마다 연산 두 번만 한다. 13장에서 보았듯이 이런 계산은 메모리에 묶여 프로세서의 연산 능력을 쓰지 못한다. 선형대수의 기본 연산을 \term{BLAS}{basic linear algebra subprograms}라는 표준 인터페이스로 묶고 그 층을 셋으로 나누면 이 문제가 분명해진다. 1층은 벡터 연산($y\gets\alpha x+y$, 내적)으로 데이터 $O(n)$에 연산 $O(n)$이고\cite{lawson1979blas1}, 2층은 행렬--벡터 연산으로 데이터 $O(n^2)$에 연산 $O(n^2)$이며, 3층은 행렬--행렬 연산($C\gets C+AB$)으로 데이터 $O(n^2)$에 연산 $O(n^3)$이다\cite{dongarra1990blas3}. 데이터 하나당 연산의 수가 $O(n)$인 것은 3층뿐이므로, 블록화(13장)로 캐시를 충분히 재사용해 연산 장치를 쉬지 않게 할 수 있는 것도 3층뿐이다.

\subsection{Block LU Factorization}

LU 분해를 3층 연산 위주로 다시 짜려면 행렬을 폭 $b$의 블록으로 나눈다. 첫 $b$개의 열, 곧 \term{패널}{panel}을 부분 피벗으로 분해해 $P_1\begin{psmallmatrix}A_{11}\\A_{21}\end{psmallmatrix}=\begin{psmallmatrix}L_{11}\\L_{21}\end{psmallmatrix}U_{11}$을 얻고, 패널의 행 교환을 나머지 열에 적용한 뒤
\[
U_{12}=L_{11}^{-1}A_{12},\qquad A_{22}\gets A_{22}-L_{21}U_{12}
\]
를 계산한다. 첫째 식은 하삼각 행렬로 여러 우변을 푸는 3층 연산 \code{trsm}이고, 둘째 식은 행렬 곱 \code{gemm}이다. 그 뒤 갱신된 $A_{22}$에 같은 일을 재귀적으로 한다. 연산의 대부분은 \code{gemm}에 있으므로 분해 전체가 행렬 곱의 속도에 가깝게 돈다. 패널의 분해는 여전히 2층 연산이지만 그 비용은 $O(n^2b)$로 전체의 작은 부분이다. 행렬을 반으로 나누어 재귀적으로 분해하는 판은 블록의 크기를 정하지 않아도 모든 층의 캐시에 저절로 맞는다\cite{toledo1997recursive}. 같은 블록 구조는 반올림 오차의 상한에도 거의 영향을 주지 않는다\cite{higham2002accuracy}. LAPACK의 \code{getrf}가 이 블록 LU이고\cite{anderson1999lapack}, \code{getf2}는 블록 없는 \cref{alg:lu-pp}이다.

\begin{instardust}[\code{docs/design/linalg.md} §1.2, §7.5, §10.2]
stardust의 선형대수 영역은 Apple Accelerate의 LAPACK을 감싼다. \code{solve}는 \code{getrf}로 분해한 뒤 \code{getrs}로 풀고, 우변이 두 열 이하이면 \code{getrs} 대신 행 교환 뒤 2층 삼각 풀기 \code{trsv}를 $L$과 $U$에 한 번씩 부른다. 설계 문서에 따르면 Accelerate의 3층 삼각 풀기는 $n\ge100$의 한 열에서 2층 판보다 $3$--$4$배 느렸다. 측정에서 흥미로운 예외도 드러났다. Accelerate는 $n=32$--$40$의 \code{getrf}를 코어들이 함께 쓰는 행렬 가속 장치(AMX)에서 직렬로 돌려 스레드를 늘려도 빨라지지 않았으므로, stardust는 $n=32$--$46$의 LU를 블록 없는 \code{getf2}로 한다. 두 판의 결과는 반올림 수준에서 다를 수 있지만 피벗 규칙과 특이의 판정은 같다. 특이 행렬은 LAPACK이 정확한 $0$ 피벗을 보고할 때만으로 정하고, 조건수를 보고 ``거의 특이''를 판정하지는 않는다.
\end{instardust}

\section{Determinants}
\label{sec:determinant}

LU 분해는 행렬식도 준다. $\det(PA)=\det L\det U=\prod_ku_{kk}$이고 $\det P=\pm1$은 행 교환 횟수의 홀짝으로 정해지므로, $\det A=\pm\prod_ku_{kk}$이다. 문제는 곱이 쉽게 넘치거나 아래로 넘친다는 것이다. $1000\times1000$ 행렬의 대각 원소가 모두 $3$ 근처이면 곱은 $3^{1000}\approx10^{477}$ 근처로 binary64의 범위를 넘고, $1/3$ 근처이면 아래로 넘친다. 그래서 행렬식의 로그 $\log|\det A|=\sum_k\log|u_{kk}|$와 부호를 따로 구하는 것이 표준이다. 곱을 직접 구할 때도 각 인수를 $m\cdot2^e$로 나누어 유효숫자와 지수를 따로 모으되, 곱할 때마다 누적한 유효숫자를 다시 유효숫자와 지수로 나누어 두면 유효숫자의 곱은 늘 $[1/2,1)$에 머물고 지수는 정수로 정확히 더해지므로 넘침 없이 계산할 수 있다.

\begin{instardust}[\code{csrc/linalg/ops/lu.cpp}]
\code{det}과 \code{slogdet}은 LU 분해의 대각 곱과 피벗 교환의 부호로 계산한다. 주석에 적힌 대로 곱을 구할 때는 각 값을 큰 성분이 $[0.5,1)$에 오는 유효숫자 $m$과 지수 $e$로 나누어 따로 모은다. 특이 행렬은 오류가 아니라 행렬식 $0$, 로그 행렬식 $-\infty$로 돌려준다. 복소 행렬의 LU에서 \code{zgetrf}는 $|\operatorname{Re}|+|\operatorname{Im}|$이 가장 큰 원소를 피벗으로 고르므로, 보장되는 것은 $|l_{ij}|\le1$이 아니라 $|l_{ij}|\le\sqrt2$이고 테스트도 이 경계로 확인한다.
\end{instardust}

\section{Concluding Remarks}

가우스 소거는 가장 오래된 수치 알고리즘이면서 반올림 오차 해석이 처음으로 성공한 알고리즘이기도 하다. 계산한 분해는 조금 바뀐 행렬의 정확한 분해이고, 그 ``조금''의 크기는 성장 인수가 정한다. 성장 인수가 지수적으로 커질 수 있다는 최악의 경우와 실제로는 거의 언제나 작다는 경험 사이의 간극은 아직 완전히 메워지지 않았지만, 부분 피벗 LU는 그 간극을 감수하고 수십 년째 연립방정식의 표준 해법으로 쓰인다. 다음 장에서는 행렬이 대칭일 때 이 분해를 절반의 비용으로, 그리고 피벗 없이 안정하게 할 수 있음을 본다.

가우스 소거의 오차 해석은 Higham의 책\cite{higham2002accuracy} 9장이 가장 완전하고, 블록 알고리즘은 Golub과 Van Loan의 책\cite{golub2013matrix}과 Demmel의 교과서\cite{demmel1997}에 자세하다.

\chapterreferences
